\section{Introdução}

%%=============================================================================
%% NOTA DE TRABALHO -- LEIA ANTES DE ESCREVER
%%
%% A INTRODUÇÃO É A ÚLTIMA COISA QUE ESCREVEREMOS.
%%
%% Motivo: a introdução tem de prometer exatamente o que o trabalho entrega.
%% Enquanto a especificação, a amostra e a leitura dos resultados ainda estão
%% se estabilizando (Capítulos 3 e 4), qualquer introdução escrita agora vira
%% dívida: promete efeito causal onde há associação condicional, promete
%% horizonte longo onde a margem é trimestral, promete "impacto da IA" onde o
%% tratamento é um índice ocupacional. Escrevemos a introdução, o resumo e o
%% abstract juntos, no fim, a partir do que os Capítulos 3, 4 e 5 já disserem.
%%
%% Ordem de escrita combinada: Capítulo 3 -> Capítulo 4 -> Capítulo 5 ->
%% Capítulo 2 (fechamento da literatura) -> Capítulo 1 -> resumo/abstract.
%%=============================================================================

Esta seção será redigida ao final, depois de fechados os procedimentos
metodológicos, os resultados preliminares e as considerações. O objetivo é
garantir que a promessa feita aqui corresponda exatamente ao que o desenho
identifica e ao que os dados sustentam.

Quando for escrita, a introdução deverá cobrir, nesta ordem:

\begin{enumerate}[label=\alph*)]
  \item \textbf{contextualização}: a difusão de modelos de linguagem de grande
        porte a partir do lançamento público do ChatGPT, em 30 de novembro de
        2022, e por que esse marco é usado como referência temporal;
  \item \textbf{problema de pesquisa}: trabalhadores em ocupações mais expostas
        à inteligência artificial passaram a apresentar transições de trabalho
        diferentes (mudança de ocupação, saída do emprego, informalização,
        deslocamento no gradiente de exposição) depois desse marco?
  \item \textbf{hipóteses}: as hipóteses de trabalho que serão confrontadas com
        os desfechos do Capítulo~\ref{sec:resultados}, escritas de forma que
        possam ser rejeitadas;
  \item \textbf{objetivo geral e objetivos específicos}: medir o gradiente
        condicional de transição ocupacional por desvio-padrão de exposição e
        descrever a mobilidade entre faixas de exposição, deixando explícito que
        a validação do desenho (estudo de evento, testes de tendência
        pré-marco e bateria de sensibilidade) fica para a agenda do
        Capítulo~\ref{sec:melhorias};
  \item \textbf{justificativa}: a inexistência, no Brasil, de medida direta de
        adoção de inteligência artificial em bases representativas, e o que uma
        medida de exposição ocupacional aplicada a um painel domiciliar
        rotativo permite e não permite concluir;
  \item \textbf{delimitação}: o que o trabalho \emph{não} pretende identificar,
        antecipando a Seção~\ref{sec:limites};
  \item \textbf{estrutura do texto}: um parágrafo apresentando os cinco
        capítulos.
\end{enumerate}

% TODO: substituir todo o conteúdo acima pelo texto final da introdução.
